\begin{figure}[tbp]
\centering
\includegraphics[width=\textwidth]{main/figure_01.pdf}
\caption{\textbf{Task-derived credit separates neuron identity from resolution within a dendritic tree.} \textbf{A}, Synaptic conductance changes combine a local eligibility with delivered compartment credit. Here $\eta$ is the learning rate; $x_i$ is presynaptic activity; $R_n$ is total input resistance; $E_i-V_n$ is the synaptic driving force; and $\varepsilon_n$ is the delivered voltage-error signal. In a directed tree, exact transport multiplies somatic error $\delta_0$ by local path derivatives $\alpha$. Reciprocal cables require an adjoint solution. \textbf{B}, An illustrative neuron with six nonsomatic compartments supports an all-site broadcast, two subtree profiles, or six independent compartment profiles. The delivery matrix A has one row per compartment and one column per spatial profile; c supplies their example-dependent coefficients, and K counts profiles. These are possible field representations, not six independently supplied external errors: exact paths can transform one somatic error using state-dependent local derivatives. \textbf{C}, Neuronal coordinates are supplied either by the task-readout derivative or a fixed random output-to-neuron map, termed direct feedback alignment (DFA). The choice between neurons is crossed with the resolution used within each tree. The layer-scalar control reduces neuronal coordinates to a shared scalar as defined in Methods; the per-neuron condition broadcasts each neuron's own coordinate through its arbor. \textbf{D}, MNIST test accuracy with readout-derived coordinates in shunting and raw-additive trees. Thin lines join paired seeds; large symbols and intervals show means and retained 95\% confidence intervals (15 seeds per architecture). \textbf{E–F}, Paired accuracy differences for preserving neuron identity (E) and then adding exact path resolution (F), using readout or DFA sources on MNIST and a readout source on flattened CIFAR-10. Colors identify architectures as in D; 'additive' denotes the raw-additive comparator. CIFAR-10 uses 20 seeds. Existing confidence intervals are reproduced without recomputation; pp denotes percentage points. MNIST's raw-additive path increment is small and positive, whereas CIFAR-10's is negative, bounding the claim to the predominance of neuron identity in these image tasks.}
\label{fig:framework}
\label{fig:feedback}
\end{figure}
